\section{Introduction}

This report analyses the paper by Binstok, Gould and Kaplan (2024) and addresses its research question, methods and results. We focus on the replication of the ALLBUS analysis and its main findings, and we give a short overview of our replication, our critique and our extension of the paper.

\subsection*{Research Problem}

The paper addresses the question of whether personal prejudices shape political views. The existing literature on immigration, race and the War on Terror finds a correlation between prejudice and political views, but provides little causal evidence. A central problem is that prejudice is difficult to measure in surveys: social desirability creates a bias, so that respondents hide their prejudices and their answers do not reflect their true views.

The paper applies this question to a prominent public debate: whether criticism of Israel is motivated by antisemitism. Critics of Israel argue that the accusation of antisemitism is used to silence legitimate criticism, while Jewish advocacy groups argue that parts of this criticism draw on antisemitic stereotypes. 
Empirical evidence on the link between antisemitism and views on Israel is limited and not causal; Kaplan and Small (2006), for example, only find a correlation.

\subsection*{Research Question}

The central research question is therefore whether antisemitic beliefs cause more negative views about Israel. 

\subsection*{Our Replication and Extension}



\par\textbf{Empirical Approach}\par

\begin{itemize}
\item
  The paper combines two independent empirical strategies, each relying on its own instrumental variable.\footnote{(Binstok, Gould, and Kaplan 2024, p. 2332, l. 26--28; p. 2335, l. 11--14)}
\item
  \textbf{Strategy 1, Twitter (United Kingdom, 2013--2022):} a panel of 355 districts observed over 3,531 days, estimated by OLS and IV with district and time fixed effects.\footnote{(Binstok, Gould, and Kaplan 2024, p. 2333, l. 6--20; p. 2335, l. 16--25)}

  \begin{itemize}
  \item
    Instrument: online chatter about Tottenham Hotspur, a football club widely perceived as a ``Jewish'' team but unrelated to Israel.\footnote{(Binstok, Gould, and Kaplan 2024, pp. 2332--2333, l. 29--35 and l. 1--8)}
  \end{itemize}
\item
  \textbf{Strategy 2, German General Social Survey (ALLBUS 2016):} a cross-section of individuals, estimated by OLS with successively larger sets of controls and subsequently by IV.\footnote{(Binstok, Gould, and Kaplan 2024, pp. 2333--2334, l. 21--34 and l. 1--2)}

  \begin{itemize}
  \item
    Instrument: being born between 1930 and 1945, i.e. exposure to Nazi indoctrination during childhood (following Voigtländer and Voth 2015), compared with respondents born 1946--1960.\footnote{(Binstok, Gould, and Kaplan 2024, p. 2334, l. 3--9)}
  \end{itemize}
\item
  Robustness: a 2004 survey of the Anti-Defamation League (ADL) covering ten European countries, again with birth cohort as instrument.\footnote{(Binstok, Gould, and Kaplan 2024, p. 2334, l. 12--13; p. 2365, l. 39--41)}
\end{itemize}

\par\textbf{Main Findings}\par

\begin{itemize}
\item
  \textbf{Twitter:} negative tweets about Jews are strongly correlated with negative tweets about Israel at the district-day level.\footnote{(Binstok, Gould, and Kaplan 2024, p. 2333, l. 12--13)}

  \begin{itemize}
  \item
    The OLS coefficient is 0.297, the IV coefficient 0.748, i.e. more than twice as large, with a first-stage F-statistic above 25.\footnote{(Binstok, Gould, and Kaplan 2024, p. 2340, l. 7--12)}
  \item
    In a placebo test, negative tweets about Jews have no effect on negative tweets about China.\footnote{(Binstok, Gould, and Kaplan 2024, p. 2333, l. 18--20; p. 2340, l. 35--38)}
  \item
    The effect on positive tweets about Israel is small.\footnote{(Binstok, Gould, and Kaplan 2024, p. 2340, l. 30--34)}
  \end{itemize}
\item
  \textbf{ALLBUS, OLS:} agreement that ``Jews have too much influence'' is positively associated with agreement that ``Israel is unjust to take Palestinian land''; the coefficient rises from 0.118 to 0.162 as controls are added (Table 7, upper panel, p. 2353).\footnote{(Binstok, Gould, and Kaplan 2024, p. 2352, l. 11--15 and l. 34--37)}
\item
  \textbf{ALLBUS, IV:} the estimated effects are considerably larger than the OLS estimates and also increase with the set of controls.\footnote{(Binstok, Gould, and Kaplan 2024, p. 2334, l. 9--12)}

  \begin{itemize}
  \item
    In the main specification, the coefficient is 0.850 with a standard error of 0.339 (Table 10, column 4, p. 2359; repeated in Table 11).\footnote{(Binstok, Gould, and Kaplan 2024, p. 2362, l. 12--13)}
  \item
    In terms of magnitude, an increase of one standard deviation in antisemitism (factor measure) raises anti-Israel views by about 1.10 standard deviations under IV, compared with 0.18 under OLS.\footnote{(Binstok, Gould, and Kaplan 2024, p. 2364, l. 18--27)}
  \end{itemize}
\item
  The authors attribute the gap between IV and OLS to under-reporting: respondents conceal antisemitic views (measurement error) but express them through their views on Israel.\footnote{(Binstok, Gould, and Kaplan 2024, p. 2364, l. 27--35; p. 2370, l. 34--38)}
\item
  The authors stress an important qualification: about half of the respondents with the most negative views of Israel show no sign of antisemitism.\footnote{(Binstok, Gould, and Kaplan 2024, p. 2351, l. 8--20; p. 2365, l. 35--38)}
\item
  In a counterfactual exercise, capping antisemitism at level 3 lowers the mean anti-Israel view by 0.76, i.e. by roughly half a standard deviation.\footnote{(Binstok, Gould, and Kaplan 2024, p. 2365, l. 12--19)}
\item
  The ADL data yield a similar pattern, with IV estimates exceeding OLS estimates and rising with the set of controls.\footnote{(Binstok, Gould, and Kaplan 2024, p. 2370, l. 23--29)}
\end{itemize}

\par\textbf{Replication and Extension}\par

\begin{itemize}
\item
  \textbf{Replication:} the analysis focuses on the ALLBUS results, namely Table 6 (column 1), Table 7, the first stage in Table 9 (third panel, column 1), Table 10 and Figure 7. The reasons for this selection are set out in Chapter 3.
\item
  \textbf{Status:} all selected estimates have been reproduced exactly in R from the raw ALLBUS file, including the main estimate of 0.850 (0.339).
\item
  \textbf{Extension:} still to be specified. A natural candidate is the strength of the first stage, which falls to about 10 in the main specification (Section 3.4).
\end{itemize}

